\documentclass[12pt]{article}
\usepackage[utf8]{inputenc}
\usepackage[T1]{fontenc}
\usepackage{hyperref}
\title{Tomoip Language Resources for K-12 Curriculum Development}
\author{Compiled by Claude Sonnet 5 (claude-sonnet-5) for Andr\'as Kornai}
\date{}
\begin{document}
\maketitle

\section{Introduction}
Tomoip (ISO 639-3 tqp) is an Oceanic language of the Austronesian family, spoken on New Britain in Papua New Guinea. The search results report differing speaker numbers: about 2,000 in one summary and 700 in an older figure on Wikipedia. The materials found are academic. No educational materials were found.

\section{Primary Educational Resources}
No primers, textbooks or school programmes were found in the search results.

\section{Linguistic and Reference Materials}
\begin{itemize}
\item Tomoip language article on Wikipedia: \url{https://en.wikipedia.org/wiki/Tomoip_language}
\item ``Tomoip phonetics and phonology'' (NZ Linguistic Society): \url{https://nzlingsoc.org/journal_article/tomoip-phonetics-and-phonology/}
\item A Project MUSE article on Tomoip, apparently the 2025 historical phonology in \emph{Oceanic Linguistics}; I did not open it to confirm: \url{https://muse.jhu.edu/article/974305}
\item Tomoip numerals (MPI-EVA Channumerals database): \url{https://lingweb.eva.mpg.de/channumerals/Tomoip.htm}
\item Papers by Russell Barlow: \url{https://eva-mpg.academia.edu/RussellBarlow}
\end{itemize}

\section{Religious and Community Resources}
No Tomoip scripture or community organisations were found in the search results.

\section{Digital and Audio Resources}
\begin{itemize}
\item Wikitongues page for tqp: \url{https://wikitongues.org/languages/tqp/}
\end{itemize}
One search summary says the Malcolm Ross collection at PARADISEC holds Tomoip materials. The results gave no URL for it, so it is not listed as a resource.

\section{Curriculum Development Context}
Documentation is academic and very limited. A curriculum would have to start from new work with the community, using the phonology descriptions as a basis for an orthography.
\end{document}